%%%%%%%%%%%%%%%%%%%%%%%%%%%%%%%%%%%%%%%%%%
% a0poster Landscape Poster
% LaTeX Template
% Version 1.0 (22/06/13)
%
% The a0poster class was created by:
% Gerlinde Kettl and Matthias Weiser (tex@kettl.de)
% 
% This template has been downloaded from:
% http://www.LaTeXTemplates.com
%
% License:
% CC BY-NC-SA 3.0 (http://creativecommons.org/licenses/by-nc-sa/3.0/)
%
%%%%%%%%%%%%%%%%%%%%%%%%%%%%%%%%%%%%%%%%%

%----------------------------------------------------------------------------------------
%	PACKAGES AND OTHER DOCUMENT CONFIGURATIONS
%----------------------------------------------------------------------------------------

\documentclass[a0,landscape]{a0poster}

\usepackage{multicol} % This is so we can have multiple columns of text side-by-side
\columnsep=100pt % This is the amount of white space between the columns in the poster
\columnseprule=3pt % This is the thickness of the black line between the columns in the poster

\usepackage[svgnames]{xcolor} % Specify colors by their 'svgnames', for a full list of all colors available see here: http://www.latextemplates.com/svgnames-colors

%\usepackage{times} % Use the times font
\usepackage{palatino} % Uncomment to use the Palatino font

\usepackage{graphicx} % Required for including images
\graphicspath{{figures/}} % Location of the graphics files
\usepackage{booktabs} % Top and bottom rules for table
\usepackage[font=small,labelfont=bf]{caption} % Required for specifying captions to tables and figures
\usepackage{amsfonts, amsmath, amsthm, amssymb} % For math fonts, symbols and environments
\usepackage{wrapfig} % Allows wrapping text around tables and figures

\makeatletter
\newcommand\thefontsize[1]{{#1 The current font size is: \f@size pt\par}}
\makeatother

\makeatletter
\@ifdefinable{org@thebibliography}{%
  \let\org@thebibliography\thebibliography
  \def\thebibliography#1{%
    \@ifnextchar\par{%
      \par@thebibliography{#1}%
    }{%
      \check@thebibliography{#1}{}%
    }%
  }%
  \newcommand{\par@thebibliography}[2]{%
    \check@thebibliography{#1}{\par}%
  }%
  \newcommand{\check@thebibliography}[2]{%
    \@ifnextchar\end{%
      \let\endthebibliography\relax
    }{%
      \org@thebibliography{#1}#2%
    }%
  }%
}
\makeatother

\begin{document}

%----------------------------------------------------------------------------------------
%	POSTER HEADER 
%----------------------------------------------------------------------------------------

% The header is divided into three boxes:
% The first is 55% wide and houses the title, subtitle, names and university/organization
% The second is 25% wide and houses contact information
% The third is 19% wide and houses a logo for your university/organization or a photo of you
% The widths of these boxes can be easily edited to accommodate your content as you see fit

\begin{minipage}[b]{0.55\linewidth}
\veryHuge \color{NavyBlue} \textbf{Packing Heat: Guns in Public Spaces} \color{Black}\\ % Title
\Huge\textit{Gun Ownership Identity and Public Safety}\\[1cm] % Subtitle
\huge \textbf{Patrick Gauding}\\ % Author(s)
\huge University of Kansas\\ % University/organization
\end{minipage}
%
\begin{minipage}[b]{0.25\linewidth}
\color{DarkSlateGray}\Large \textbf{Contact Information:}\\
Blake Hall, Room 504\\ % Address
University of Kansas\\
1541 Lilac Ln., Lawrence KS 66044\\\\
Phone: 785.864.3523 \\ % Phone number
Email: \texttt{patrick.gauding@ku.edu}\\ % Email address
\end{minipage}
%
\begin{minipage}[b]{0.19\linewidth}
\includegraphics[width=16cm]{logo.png} % Logo or a photo of you, adjust its dimensions here
\end{minipage}

\vspace{1cm} % A bit of extra whitespace between the header and poster content

%----------------------------------------------------------------------------------------

\begin{multicols}{4} % This is how many columns your poster will be broken into, a poster with many figures may benefit from less columns whereas a text-heavy poster benefits from more

%----------------------------------------------------------------------------------------
%	ABSTRACT
%----------------------------------------------------------------------------------------

\color{Navy} % Navy color for the abstract
\huge
\begin{abstract}
\normalsize
Americans increasingly percieve that guns make individuals and society safer, but the public remains polarized at this point. Given recent expansions of concealed carry laws and 'stand your ground' policies, it is clear that public attitudes about guns and safety have real implications for public policy. This research seeks to understand the social and psychological roots of divergent attitudes on this issue. I hypothesize that social and political identities motivate beliefs about guns and safety. These identities lead their adherents to engage in a form of motivated reasoning that results in preferences and beliefs that protect self-interests. I test these hypotheses through analysis of a recent national survey of American adults. These results suggest that gun-owners, men, Republicans and conservatives are the most likely to believe that guns make us safer, which conforms to the self-interests of these groups.

\end{abstract}

%----------------------------------------------------------------------------------------
%	INTRODUCTION
%----------------------------------------------------------------------------------------

\color{SaddleBrown} % SaddleBrown color for the introduction

\huge
\section*{Research Question}

\LARGE How are attitudes towards guns in public spaces affected by gun \\ownership?

\color{DarkSlateGray} % DarkSlateGray color for the rest of the content

\section*{Literature}
\Large
\begin{itemize}
	\item[] Motivated reasoning - Kunda (1990), Lodge and Taber (2013)
	\item[] Symbolic politics of guns - Wolpert and Gimpel (1998), Celinska (2007)
	\item[] Gun ownership as an identity - Carlson (2015), Joslyn et al. (2017), Joslyn and Haider-Markel (2017)
\end{itemize}

\section*{Theory and Hypotheses}

$H_1$: Among all respondents, attitudes towards guns in public spaces and safety should be affected by gunownership, gender, ideology and party identity.\\
\\
$H_2$: Among gunowners, attitudes towards guns in public spaces and safety should be affected by strength of gun ownership as an identity.


\section*{Data and Methods}
\large
\begin{itemize}
	\item[-] National survey of American adults (\emph{N} = 2088) conducted in summer 2017
	\item[-] Oversample of gunowners; $N_{gunowners}$ = 855
	\item[-] Due to censored/truncated nature of additive index DV, analysis uses Tobit model
\end{itemize}

%------------------------------------------------

\subsection*{Variables}
\large
\begin{itemize}
	\item[-] DV: Guns and Public Safety Index (additive, coded 5-25, Cronbach's $\alpha$ = 0.81)
	\item[-] IV1: Gun in the Home (Dichotomous, coded 1 if gun present)
	\item[-] IV2: Gun Ownership Identity (additive, coded 4-28, Cronbach's $\alpha$ = 0.76)
	\item[-] Controls: Ideology (1-7), Age, Education (1-7), Income, White (dummy), Republican (dummy), Democrat (dummy)
	
\end{itemize}

%----------------------------------------------------------------------------------------
%	RESULTS 
%----------------------------------------------------------------------------------------

\section*{Results}
\large
\begin{center}\vspace{1cm}
\begin{tabular}{*{4}{l}}

     &  Model 1	&  Model 2 & Model 3\tabularnewline
   & Estimate & Estimate & Estimate \tabularnewline
    & (S.E.) & (S.E.)& (S.E.)\tabularnewline
 \hline
 \hline
  Gun in the Home    		& 5.179*** 	& 4.258*** 	& \tabularnewline
     										&  (0.209) 	&  (0.190) 		& \tabularnewline
  Gun Identity    			& 				 	&  						& 0.171***\tabularnewline
     										&  					&   					& (0.021)\tabularnewline
  Ideology    					&  					& 0.964*** 		& 0.636***\tabularnewline
     										&   				&  (0.068) 		& (0.090)\tabularnewline
  Gender    						&  					& 0.777*** 	  & -0.257\tabularnewline
     										&   				&  (0.185)  	& (0.257)\tabularnewline
  Age     							&  					& -0.011 		& 0.022*\tabularnewline
     										&   				&  (0.007)  	& (0.009)\tabularnewline
  Education     				&  				 	& -0.253***		& -0.114\tabularnewline
    										&  		 		 	&  (0.070)  	& (0.098)\tabularnewline
  Income    						&  					& -0.118**  	& 0.073\tabularnewline
     										&   				&  (0.045)  	& (0.066)\tabularnewline
	White									&						& -0.052			&	0.503\tabularnewline
												&						& (0.235)			&	(0.357)	\tabularnewline
  Republican    				&  					& 0.701*** 		& -0.181\tabularnewline
     										&   				&  (0.248)  	& (0.313)\tabularnewline
  Democrat     					&     			& -1.305*** 	& -0.885*\tabularnewline
     										&   				&  (0.236)  	& (0.364)\tabularnewline
  Constant	    				&	13.348*** &	17.047*** 	& 12.950***\tabularnewline
     										&  (0.147) 	&  (1.269) 		& (1.854)\tabularnewline
 \hline
  Var.					& 21.001		&		15.878						& 11.884\\
  	 						& (0.711) 	&  (0.537)						& (0.624)\\
\hline
 N 	 						& 1945 			& 1945  				& 807 \tabularnewline
 LR $\chi^2$    & 537.10   	& 1065.55   		& 224.28 \tabularnewline
 Prob $>$ $\chi^2$ & 0.000  & 0.000  				& 0.000  \tabularnewline
 Pseudo $R^2$   & 0.047   	& 0.092   			& 0.052 \tabularnewline
 Log-likelihood & -5508.724 & -5244.502   	& -2064.103 \tabularnewline
 \hline
\hline
 
 \multicolumn{3}{c}{${*  p}\le 0.05$${*\!\!*  p}\le 0.01$${*\!\!*\!\!*  p}\le 0.001$}\tabularnewline
 \end{tabular}
%\captionof{table}{\color{Green} Table 1}
\end{center}\vspace{1cm}

%----------------------------------------------------------------------------------------
%	CONCLUSIONS
%----------------------------------------------------------------------------------------

\color{SaddleBrown} % SaddleBrown color for the conclusions to make them stand out

\section*{Conclusions}
\large
\begin{itemize}
\item Among all respondents, support for guns in public spaces being postive for safety are affected by the presence of a gun in the home, ideology, gender, education, income, and party identity, consistent with previous research.
\item Among gunowners, gun identity, ideology, age and being a Democrat affect these attitudes. This is consistent with the hypothesized expectation that gun identity exerts a postive effect on attitudes towards guns and public safety.
\item Ideology does exert a nearly 0.5 unit higher impact on the gun-safety index than identity, which is not particularly surprising.
\item However, the strength of the gun identity effect (p $<$ 0.001) points to its importance in unwraveling differing political preferences between non-gunowners and gunowners.
\end{itemize}

\color{DarkSlateGray} % Set the color back to DarkSlateGray for the rest of the content

%----------------------------------------------------------------------------------------
%	FORTHCOMING RESEARCH
%----------------------------------------------------------------------------------------

 %----------------------------------------------------------------------------------------
%	REFERENCES
%----------------------------------------------------------------------------------------

\subsection*{References}
\small
\begin{itemize}
\item Carlson, J. (2015). \textit{Citizen-protectors: The everyday politics of guns in an age of decline.} Oxford University Press.
\item Celinska, K. (2007). "Individualism and collectivism in America: The Case of Gun Ownership and Attitudes toward Gun Control." \textit{Sociological Perspectives} 50(2): 229-247.
\item Joslyn, M. R. and Haider-Markel, D. P. (2017). "Gun ownership and self-serving attributions for mass shooting tragedies." \textit{Social Science Quarterly} 98(2):429–442.
\item Joslyn, M. R., Haider-Markel, D. P., Baggs, M., and Bilbo, A. (2017). "Emerging political identities? gun ownership and voting in presidential elections." \textit{Social Science Quarterly} 98(2):382–396.
\item Kunda, Z. (1990). "The case for motivated reasoning." \textit{Psychological bulletin}, 108(3):480.
\item Lodge, M. and Taber, C. S. (2013). \textit{The rationalizing voter.} Cambridge University Press.
\end{itemize}

%----------------------------------------------------------------------------------------
%	ACKNOWLEDGEMENTS
%----------------------------------------------------------------------------------------

\section*{Acknowledgements}

I thank Clayton Webb for his help and instruction in this class, as well as Donald P. Haider-Markel for data and support on this project.

%----------------------------------------------------------------------------------------

\end{multicols}
\end{document}